% Figure: p-value visualised as a tail area under the null distribution.
\begin{figure}[htbp]
\centering
\begin{tikzpicture}
\begin{axis}[
    width=11cm, height=5.5cm,
    xmin=-4, xmax=4,
    ymin=0, ymax=0.45,
    axis lines=left,
    xlabel={Test statistic $T$},
    ylabel={Null density $p_{H_0}(t)$},
    samples=200,
    domain=-4:4,
    no markers,
    every axis plot/.append style={very thick},
    xtick={-1.96, 0, 1.96, 2.4},
    xticklabels={$-z_{1-\alpha/2}$, $0$, $z_{1-\alpha/2}$, $T(x)$},
    ytick=\empty,
    clip=false,
]
\addplot[engblue]
  {1/sqrt(2*pi)*exp(-x*x/2)};

\addplot[engred!50, fill=engred!30, fill opacity=0.7,
         domain=2.4:4, samples=80]
  {1/sqrt(2*pi)*exp(-x*x/2)} \closedcycle;

\addplot[engred!50, fill=engred!30, fill opacity=0.7,
         domain=-4:-2.4, samples=80]
  {1/sqrt(2*pi)*exp(-x*x/2)} \closedcycle;

\draw[engblack, dashed] (axis cs:2.4, 0) -- (axis cs:2.4, 0.30);
\node[engred!70!black, anchor=west, font=\small] at (axis cs:2.5, 0.18)
  {$p/2$};

\draw[engblack, dashed] (axis cs:1.96, 0) -- (axis cs:1.96, 0.18);
\draw[engblack, dashed] (axis cs:-1.96, 0) -- (axis cs:-1.96, 0.18);
\end{axis}
\end{tikzpicture}
\caption[The two-sided $p$-value as a tail area]{The two-sided
$p$-value of a realised statistic $T(x)$ is the probability, under
the null distribution, of producing a statistic at least as extreme
as $T(x)$. Shaded regions are the two tails beyond $\pm \abs*{T(x)}$.
The critical values $\pm z_{1-\alpha/2}$ delimit the rejection
region at significance level $\alpha$. A $p$-value is not the
probability that the null is true and not the probability that the
alternative is true; it is a property of the test procedure given
the null.}
\label{fig:vol02:ch14:pvalue}
\end{figure}
